For simplicity, our focus thus far has been on classifiers that are perfectly
calibrated. Here, we introduce an approximate notion of calibration, which we
will use in subsequent proofs.
%
\begin{definition}
  \label{def:calibration}
  The \emph{calibration gap} $\calib{h_t}$ of a classifier $h_t$
  with respect to a group $G_t$ is
  %
  \begin{align}
    \calib{h_t}
    &= \int_{0}^{1} \Bigl\lvert \Pr_{(\x, y) \sim G_t} \bigl[ y\!=\!1 \mid h(\x)\!=\!p
  \bigr] - p \Bigr\rvert \Pr_{(\x, y) \sim G_t} \bigl[ h(\x) = p \bigr] dp.
      \label{eqn:approx-cal}
      %\nonumber
  \end{align}
  %
  Thus, a classifier $h_t$ is \emph{perfectly calibrated} if $\calib{h_t} = 0$.
\end{definition}
%
A majority of this supplementary material is devoted to proving approximate
versions of our major findings. In all cases, our results degrade smoothly as
the calibration condition is relaxed.
In addition, we also provide extended details on the experiments run in this paper.

Note that we will use the notational abuse
$\Pr_{G_t}$ and $\ev_{G_t}$ in place of $\Pr_{(\x, y) \sim G_t}$
and $\ev_{(\x, y) \sim G_t}$.

\section{Linearity of Calibrated Classifiers}
\label{app:linearity}

  In \autoref{sec:setup}, we claim that the set of all calibrated classifiers
  $\mathcal{H}_t^*$ for group $G_t$ form a line in the generalized false-positive/false-negative
  plane. The following proof of this claim is adapted from \cite{kleinberg2016inherent}.
  %
\begin{lemma}
  For a group $G_t$, if a classifier $h_t$ has $\epsilon(h_t) \le
  \delta_{cal}$, then
  \begin{align*}
    \bigl| \mu_t \fn{h_t} - \left( 1 - \mu_t \right) \fp{h_t} \bigr|
    \leq 2 \delta_{cal}.
  \end{align*}
  where $\fp{h_t}$ and $\fn{h_t}$ are the generalized false-positive and false-negative
  and $\mu_t$ is the base rate of group $G_t$.
  \label{lem:approx-linear-cal}
\end{lemma}
\begin{proof}
  First, note that
  \begin{align*}
    \fp{h_t} &= \ev_{G_t} \bigl[ h_t(\x) \mid y \! = \! 0 \bigr]
    \\&= \int_{0}^1 p \: \Pr_{G_t} \left[ h_t(\x) \! = \! p \mid y \! = \! 0 \right] dp
    \\&= \int_{0}^1 p \: \frac{1 - \Pr_{G_t} \left[ y \! = \! 1 \mid h_t(\x) \! = \! p \right]}{1 - \Pr_{G_t} \left[ y \! = \! 1 \right]}
     \Pr_{G_t} \left[ h_t(\x) \! = \! p \right] dp
    \\&= \frac{1}{1-\mu_t} \int_{0}^1 p \: (1 - \Pr_{G_t} \left[ y \! = \! 1 \mid h_t(\x) \! = \! p \right])
     \Pr_{G_t} \left[ h_t(\x) \! = \! p \right] dp
    \numberthis \label{eqn:cfp-int}
    %\\&= \frac{1}{1-\mu_t} (\ev_{G_t}[h_t(\x)] -  \int_{0}^1 p \Pr_{G_t} \left[ y \! = \! 1 \mid h_t(\x) \! = \! p \right]
    %\\&\phantom{=} \: \: \times \Pr_{G_t} \left[ h_t(\x) \! = \! p \right] dp)
  \end{align*}
  Next, observe that
  \begin{align*}
    \int_0^1 & p \cdot \Pr_{G_t} \left[ y \! = \! 1 \mid h_t(\x) \! = \! p \right] \cdot
    \Pr_{G_t} \left[ h_t(\x) \! = \! p \right] dp
    \\&= \int_0^1 p (p + \Pr_{G_t} \left[ y \! = \! 1 \mid h_t(\x) \! = \! p \right] -
    p)
    \Pr_{G_t} \left[ h_t(\x) \! = \! p \right] dp
    \\&\le \int_0^1 \left(p^2 + |\Pr_{G_t} \left[ y \! = \! 1 \mid h_t(\x) \! = \! p \right] -
  p|\right)
  \Pr_{G_t} \left[ h_t(\x) \! = \! p \right] dp
    \\&\le \ev_{G_t}[h_t(\x)^2] + \delta_{cal}
  \end{align*}
  Similarly,
  \begin{align*}
    \int_0^1 & p \cdot \Pr_{G_t} \left[ y \! = \! 1 \mid h_t(\x) \! = \! p \right]
    \Pr_{G_t} \left[ h_t(\x) \! = \! p \right] dp
    \\&\ge \int_0^1 \left(p^2 - |\Pr_{G_t} \left[ y \! = \! 1 \mid h_t(\x) \! = \! p \right] -
  p|\right)
  \Pr_{G_t} \left[ h_t(\x) \! = \! p \right] dp
    \\&\ge \ev_{G_t}[h_t(\x)^2] - \delta_{cal}
  \end{align*}
  Plugging these into \eqref{eqn:cfp-int}, we have
  \begin{align*}
    \frac{1}{1-\mu_t} &(\ev_{G_t}[h_t(\x)] - \ev_{G_t}[h_t(\x)^2] -
    \delta_{cal}) \le \fp{h_t}
    \\&\le \frac{1}{1-\mu_t} (\ev_{G_t}[h_t(\x)] - \ev_{G_t}[h_t(\x)^2] +
    \delta_{cal}) \numberthis \label{eqn:fpbound}
  \end{align*}
  We follow a similar procedure for $\fn{h_t}$:
  \begin{align*}
    \fn{h_t} &= \ev_{G_t} \bigl[ 1 - h_t(\x) \mid y \! = \!0 \bigr]
    \\&= \int_{0}^1 \left( 1 - p \right) \: \Pr_{G_t} \left[ h_t(\x) \! = \! p \mid y \! = \! 0 \right] dp
    \\&= \int_{0}^1 \left( 1 - p \right) \: \frac{\Pr_{G_t} \left[ y \! = \! 1 \mid h_t(\x) \! = \! p \right]}{\Pr_{G_t} \left[ y \! = \! 1 \right]}
     \Pr_{G_t} \left[ h_t(\x) \! = \! p \right] dp.
    \\&= \frac{1}{\mu_t} \int_{0}^1 \left( 1 - p \right) (\Pr_{G_t} \left[ y \! = \! 1 \mid h_t(\x) \! = \! p \right])
    \Pr_{G_t} \left[ h_t(\x) \! = \! p \right] dp.
  \end{align*}
  We use the fact that
  \begin{align*}
    (1 - p)& (\Pr_{G_t} \left[ y \! = \! 1 \mid h_t(\x) \! = \! p \right])
    \\&=(1 - p) (p + \Pr_{G_t} \left[ y \! = \! 1 \mid h_t(\x) \! = \! p \right] - p)
    \\&\le
    p(1 - p) + |\Pr_{G_t} \left[ y \! = \! 1 \mid h_t(\x) \! = \! p \right] - p|
  \end{align*}
  and
  \begin{align*}
    (1 - p)& (\Pr_{G_t} \left[ y \! = \! 1 \mid h_t(\x) \! = \! p \right])
    \\&\ge
    p(1 - p) - |\Pr_{G_t} \left[ y \! = \! 1 \mid h_t(\x) \! = \! p \right] - p|
  \end{align*}
  to get
  \begin{align*}
    \frac{1}{\mu_t} &(\ev_{G_t}[h_t(\x)] - \ev_{G_t}[h_t(\x)^2] -
    \delta_{cal}) \le \fn{h_t}
    \\&\le \frac{1}{\mu_t} (\ev_{G_t}[h_t(\x)] - \ev_{G_t}[h_t(\x)^2] +
    \delta_{cal}) \numberthis \label{eqn:fnbound}
  \end{align*}
  Combining \eqref{eqn:fpbound} and \eqref{eqn:fnbound}, we have
  \begin{align*}
    \fn{h_t} &\le
    \frac{1}{\mu_t} (\ev_{G_t}[h_t(\x)] - \ev_{G_t}[h_t(\x)^2] +
    \delta_{cal})
    \\&= \frac{1}{\mu_t} (\ev_{G_t}[h_t(\x)] - \ev_{G_t}[h_t(\x)^2] -
    \delta_{cal} + 2\delta_{cal})
    \\&\le \frac{1}{\mu_t} ( \left(1 - \mu_t\right) \fp{h_t} + 2\delta_{cal})
    \\&= \frac{1 - \mu_t}{\mu_t} \fp{h_t} + \frac{2\delta_{cal}}{\mu_t}
  \end{align*}
  We can get a similar lower bound for \fn{h_t} as
  \begin{align*}
    \fn{h_t} &\ge
    \frac{1}{\mu_t} (\ev_{G_t}[h_t(\x)] - \ev_{G_t}[h_t(\x)^2] -
    \delta_{cal})
    \\&\ge \frac{1}{\mu_t} (\left( 1 - \mu_t \right) \fp{h_t} - 2\delta_{cal})
    \\&= \frac{1 - \mu_t}{\mu_t} \fp{h_t} - \frac{2\delta_{cal}}{\mu_t}
  \end{align*}
  Multiplying these inequalities by $\mu_t$ completes this proof.
\end{proof}
%
\begin{corollary}
  Let $\mathcal H_t$ be the set of perfectly calibrated classifiers
  for group $G_t$ --- i.e. for any $h^*_t \in \mathcal H_T$, we have
  $\calib{h^*_t} = 0$.
  The generalized false-positive and false-negative rates of
  $h^*_t$ are given by
  %
  \begin{align}
    \fp{h^*_t} &= \frac{1}{1 - \mu_t} \left( \ev_{G_t}[h_t(\x)] - \ev_{G_t}[h_t(\x)^2] \right)
      \label{eqn:fp_exact} \\
    \fn{h^*_t} &= \frac{1}{\mu_t} \left( \ev_{G_t}[h_t(\x)] - \ev_{G_t}[h_t(\x)^2] \right)
      \label{eqn:fn_exact}
  \end{align}
\end{corollary}
%
\begin{proof}
  This is a direct consequence of \eqref{eqn:fpbound} and \eqref{eqn:fnbound}.
\end{proof}
%
\begin{corollary}
  \label{crly:linearity}
  For a group $G_t$, any perfectly calibrated classifier $h^*_t$ satisfies
  %
  \begin{equation}
    \fn{h^*_t} = \frac{1 - \mu_t}{\mu_t} \fp{h_t}.
    \label{eqn:linear-calib-sup}
  \end{equation}
\end{corollary}
%
In other words, all perfectly calibrated classifiers $h^*_t \in \mathcal H_t$ for group $G_t$
lie on a line in the generalized false-positive/false-negative plane, where the slope of the
line is uniquely determined by the group's base-rate $\br_t$.

\section{Cost Functions}
\label{app:cost-funcs}

We will prove a few claims about cost functions
$\fairconstname$ of the form given by \eqref{eqn:cost_general} ---
i.e.
%
\[ \fairconst{h_t}{t} = a_t \fp{h_t} + b_t \fn{h_t} \]
for some non-negative constants $a_t$ and $b_t$.
First, we show that $h^{\br_t}$ is the calibrated classifier that maximizes
$\fairconstname$.

%\begin{lemma}
  %Let $h_t$ and $h'^t$ be two classifiers for group $G_t$.
  %$\var_{G_t} [ h_t(\x) ] \leq \var_{G_t} [ h'_t(\x) ]$ if
  %and only if
  %$\fairconst{h_t}{t} \geq \fairconst{h'_t}{t}$ for any cost function $\fairconstname$.
%\end{lemma}

\begin{lemma}
  For any cost function $\fairconstname$ that follows the form of
  \eqref{eqn:cost_general}, the trivial classifier $h^{\br_t}$ is the calibrated
  classifier for $G_t$ with maximum cost.
  \label{lma:trivial}
\end{lemma}
%
\begin{proof}
  %Let $r_t = \Pr_{G_t} \left[ y = 1 \right]$ be the base rate for group $G_t$.
  %It can be seen that for any calibrated classifier that $\ev_{G_t} \left[ h(\x) \right] = r_t$.
  Again, let $\fairconstname$ be a cost function:
  $$\fairconst{h}{t} = a_t \fp{h_t} + b_t \fn{h_t}.$$
  %
  Using \eqref{eqn:fp_exact} and \eqref{eqn:fn_exact}, we have that, for every classifier $h_t$
  that is perfectly calibrated for group $G_t$,
  %
  \begin{align*}
    \fairconst{h_t}{t} &= a_t \fp{h_t} + b_t \fn{h_t}
    \\&= \left( \frac{a_t}{1 - \mu_t} + \frac{b_t}{\mu_t} \right) \left( \ev_{G_t} \left[ h_t(x) \right] - \ev_{G_t} \left[ h_t(x)^2 \right] \right)
    \\&= \left( \frac{a_t}{1 - \mu_t} + \frac{b_t}{\mu_t} \right) \left( \mu_t - \ev_{G_t} \left[ h_t(x)^2 \right] \right).
  \end{align*}
  %
  The last equation holds because $\ev_{G_t} \left[ h_t(\x) \right] = \mu_t$ for any calibrated
  classifier -- a fact which can easily be derived from \autoref{def:calibration}.

  We would like to find, $h^{\max}_t \in \mathcal{H}^*_t$, the calibrated classifier with the highest weighted cost.
  Because $\left( \frac{a_t}{1 - \mu_t} + \frac{b_t}{\mu_t} \right)$ and $\mu_t$ are non-negative constants, we have
  %
  \begin{align*}
    h^{\max}_t &= \argmax_{h \in \mathcal{H}^*_t} \Bigl[ \left( \frac{a_t}{1 - \mu_t} + \frac{b_t}{\mu_t} \right) \left( \mu_t - \ev_{G_t} \left[ h(x)^2 \right] \right) \Bigr]
    \\&= \argmax_{h \in \mathcal{H}^*_t} \Bigl[ - \ev_{G_t} \left[ h(x)^2 \right] \Bigr]
    \\&= \argmin_{h \in \mathcal{H}^*_t} \Bigl[ \ev_{G_t} \left[ h(x)^2 \right] \Bigr]
    \\&= \argmin_{h \in \mathcal{H}^*_t} \Bigl[ \ev_{G_t} \left[ h(x)^2 \right] - \mu_t^2 \Bigr]
  \end{align*}
  %
  Thus, the calibrated classifier with minimum variance will have the highest cost. This translates to
  a classifier that outputs the same probability for every sample. By the calibration constraint, this
  constant must be equal to $\mu_t$, so this classifier must be the trivial classifier $h^{\br_t}$ --- i.e.
  for all $\x$
  \[ h^{\max}_t \left( \x \right) = h^{\br_t} \left( \x \right) = \mu_t. \]
\end{proof}

Next, we show that $\fairconstname$ is linear under randomized interpolations.
\begin{lemma} \label{lma:interpolation}
  Let $\fair h_2$ be the classifier derived from \eqref{eqn:interpolation}
  with interpolation parameter $\alpha \in [0, 1]$.
  The cost of $\fair h_2$ is given by
  \[
    \fairconst{\fair h_2}{2} = (1 - \alpha) \fairconst{h_2}{2} + \alpha \fairconst{h^{\br_2}}{2}
  \]
\end{lemma}
%
\begin{proof}
  %
  The cost of $\fair h_2$ can be calculated using linearity of expectation. Let
  $B$ be a Bernoulli random variable with parameter $\alpha$.
  %
  \begin{align*}
    \fairconst{\fair h_2}{2} &= a_2 \fp{\fair h_2} + b_2 \fn{\fair h_2}
    \\&= a_2 \ev_{G_2} \left[ 1 -  \fair h_2(\x) \mid y\!=\!1 \right] + b_2 \ev_{G_2} \left[ \fair h_2(\x) \mid y\!=\!0 \right]
    %
    \\&= a_2 \ev_{B,G_2} \left[ 1 - \left[ (1-B) h_2(\x) + B h^{\br_2}(\x) \right] \mid y\!=\!1 \right]
    %\\&\phantom{=} \:
    + b_2 \ev_{B,G_2} \left[ \left[ (1-B) h_2(\x) + B h^{\br_2}(\x) \right] \mid y\!=\!0 \right]
    %
    \\&=                a_2 \ev_{B,G_2} \left[ (1-B) \left( 1 - h_2(\x) \right) \mid y\!=\!1 \right]
    %\\&\phantom{=} \:
    + a_2 \ev_{B,G_2} \left[ B \left( 1 -  h^{\br_2}(\x) \right) \mid y\!=\!1 \right]
    \\&\phantom{=} \:
    + b_2 \ev_{B,G_2} \left[ (1-B) h_2(\x) \mid y\!=\!0 \right]
    %\\&\phantom{=} \:
    + b_2 \ev_{B,G_2} \left[ B h^{\br_2}(\x) \mid y\!=\!0 \right]
    %
    \\&=                a_2 \ev_{B} \left[ 1-B \right] \ev_{G_2} \left[ 1 - h_2(\x) \mid y\!=\!1 \right]
    %\\&\phantom{=} \:
    + a_2 \ev_{B} \left[ B \right]   \ev_{G_2} \left[ 1 -  h^{\br_2}(\x) \mid y\!=\!1 \right]
    \\&\phantom{=} \:
    + b_2 \ev_{B} \left[ 1-B \right] \ev_{G_2} \left[ h_2(\x) \mid y\!=\!0 \right]
    %\\&\phantom{=} \:
    + b_2 \ev_{B} \left[ B \right]   \ev_{G_2} \left[ h^{\br_2}(\x) \mid y\!=\!0 \right]
    %
    \\&= a_2 (1-\alpha) \fp{h_2} + b_2 (1-\alpha) \fn{h_2}
    %\\&\phantom{=} \:
    + a_2 (\alpha) \fp{h^{\br_2}} + b_2 (\alpha) \fn{h^{\br_2}}
    %
    \\&= (1 - \alpha) \fairconst{h_2}{2} + \alpha \fairconst{h^{\br_2}}{2}.
  \end{align*}
\end{proof}
